\scanpage{278}
\mspage{- 108 -}\hfill IX

où les flèches verticales désignent les foncteurs évidents induits par le foncteur exact $r$, et où $\omega'$ est la flèche analogue à $\omega$ (9.3.5), avec $C$ remplacé par $C'$. On en déduit des compatibilités analogues des structures envisagées dans 9.3.8 avec le foncteur $r$, que nous laissons au soin du lecteur à expliciter.

Nous utiliserons la conséquence suivante de cette commutativité. Supposons que $\operatorname{EXTPAN}(E,F)$ ne soit pas vide, et choisissons alors un objet $X$ dans $\operatorname{EXTPAN}(E,F)$. Alors grâce à $G'\longmapsto G'\wedge X'$ (où $X'=r(X)$), la catégorie $\operatorname{EXTPAN}(E',F')$ est équivalente à la catégorie $\operatorname{EXT}(Q',P')$, donc un élément $Y'$ de $\operatorname{EXTPAN}(E',F')$ est connu à isomorphisme près quand on connaît la classe d'isomorphie de l'élément correspondant de $\operatorname{EXT}(Q',P')$, i.e. quand on connaît l'élément correspondant de $\operatorname{Ext}^1(Q',P')$, soit

\[
c_X(Y')\in\operatorname{Ext}^1(Q',P')\,.
\]

Lorsqu'on remplace $X$ par un autre objet $X_1=H\wedge X$, où $H$ est une extension de $Q$ par $P$, on voit aussitôt que l'on obtient

\[
c_{X_1}(Y')=c_X(Y)-\mathrm{cl}(H')\qquad,
\]

grâce à la formule d'associativité (9.3.7). Par suite, pour $Y'$ fixé, et $X$ variable, les $c_X(Y')$ parcourent exactement une classe

\[
c(Y')\in\operatorname{Ext}^1(Q',P')/\operatorname{Im}\operatorname{Ext}^1(Q,P)\qquad,\tag*{(9.3.13)}
\]

qui est donc canoniquement associée à l'extension panachée $Y'$ de $E'$ par $F'$. Par construction, sa nullité est nécessaire et suffisante pour que l'extension panachée $Y'$ soit isomorphe à la transformée par $r$ d'une extension panachée $X$ de $E$ par $F$.

% Typescript upside down in the archive image.
% archive sheet number: 278
% stamp: 426 (number in pencil and on label).
